% Codebook IdeoGraphCO v1.1 — fuente LaTeX (xelatex)
% Compilar:  xelatex codebook_v1.1.tex   (dos veces, por el índice)
\documentclass[11pt,letterpaper]{article}

\usepackage{geometry}
\geometry{margin=2.4cm,top=2.6cm,bottom=2.4cm}
\usepackage{fontspec}
\usepackage{xcolor}
\usepackage{booktabs}
\usepackage{longtable}
\usepackage{array}
\usepackage{framed}
\usepackage[spanish,es-noquoting,es-nolists,es-tabla]{babel}
\usepackage[hidelinks]{hyperref}

\setmainfont{Georgia}
\newfontfamily\headfont{Helvetica Neue}

\definecolor{tinte}{HTML}{F2F4F7}
\definecolor{pend}{HTML}{FFF4E0}
\definecolor{ratif}{HTML}{E9F2EC}
\definecolor{nuevo}{HTML}{E8EEF8}
\definecolor{gris}{HTML}{6B7280}
\definecolor{azul}{HTML}{1F3864}

% ---- encabezado y pie sin fancyhdr ----
\makeatletter
\def\ps@cb{%
  \def\@oddhead{\footnotesize\color{gris}\hfill Codebook — Orientación Ideológica en Noticias Políticas Colombianas}%
  \let\@evenhead\@oddhead
  \def\@oddfoot{\hfil\footnotesize\color{gris}\thepage\hfil}%
  \let\@evenfoot\@oddfoot}
\makeatother
\pagestyle{cb}

\renewcommand{\familydefault}{\rmdefault}

% ---- cajas de aviso ----
\newcommand{\caja}[3]{%
  \par\smallskip\colorlet{shadecolor}{#1}%
  \begin{shaded}\noindent{\headfont\bfseries #2}\par\smallskip\small #3\par\end{shaded}\smallskip}

% ---- metadatos repetidos al pie de cada ficha ----
\newcommand{\metaficha}[1]{%
  \par\smallskip\noindent\footnotesize\color{gris}%
  \textbf{Correspondencia V-Party:} #1\\
  \textbf{Tipo de variable:} H (Gold Set) y LLM (Silver Set) — ver §1.4.1\\
  \textbf{Unidad de codificación:} Artículo periodístico completo (ver §1.6)\\
  \textbf{Agregación entre codificadores:} Consenso (Gold Set) / Alfa de Krippendorff (Silver Set, auditoría)\\
  \textbf{Versión del codebook:} 1.1\par\normalsize\color{black}\medskip}

% ---- ejemplos ----
\newcommand{\ejemplo}[2]{%
  \par\smallskip\noindent\textbf{Ejemplo — #1:}\par\nopagebreak
  \begingroup\itshape\small\leftskip=1.2em\rightskip=0.5em\noindent #2\par\endgroup\smallskip}

\newcommand{\dimtitulo}[2]{\subsubsection*{#1 \normalfont\small\color{gris}(H,\,LLM)\ \ \texttt{#2}}%
  \addcontentsline{toc}{subsubsection}{#1}}

\setlength{\parskip}{0.45em}
\setlength{\parindent}{0pt}

\begin{document}

% ============================ PORTADA ============================
\thispagestyle{empty}
\begin{center}
\vspace*{1.2cm}
{\headfont\large UNIVERSIDAD DEL VALLE}\\[0.25em]
{\headfont\small Facultad de Ingeniería — Escuela de Ingeniería de Sistemas y Computación}\\[2.6cm]

{\headfont\fontsize{34}{38}\selectfont\color{azul}\bfseries CODEBOOK}\\[0.9em]
{\headfont\large Manual de Reglas de Anotación de Orientación Ideológica}\\[0.3em]
{\headfont\large para Noticias Políticas Colombianas}\\[1.8cm]

{\headfont\large\bfseries v1.1 — Octubre 2026}\\[0.5em]
{\small\color{gris}Revisión de calibración de la dimensión Institucionalismo\\y ratificación de las reglas de agregación}\\[1.6cm]

{\footnotesize\color{gris}\begin{minipage}{0.86\textwidth}\centering
Estructura y convenciones adaptadas del \emph{Codebook Varieties of Party Identity and Organization}
(V-Party) V2, V-Dem Institute, University of Gothenburg (Lindberg et al., 2022)
\end{minipage}}\\[2.2cm]
\end{center}

\noindent Desarrollado para el proyecto de grado:\\
\emph{“Desarrollo de clasificador multiclase para la identificación de orientación ideológica en noticias
políticas colombianas”}

\vspace{0.8em}
\noindent Kevin Steven Trujillo Riascos — Juan Esteban Camargo Chacón\\
Director: Raúl Ernesto Gutiérrez de Piñerez Reyes, Ph.D.

\newpage

% ============================ CITA Y VERSIONES ============================
\section*{Cita sugerida}
Trujillo Riascos, K. S., \& Camargo Chacón, J. E. (2026). \emph{Codebook: Manual de reglas de anotación de
orientación ideológica para noticias políticas colombianas} (v1.1) [Manual de codificación]. Escuela de
Ingeniería de Sistemas y Computación, Universidad del Valle.

Estructura documental basada en: Lindberg, S. I. et al. (2022). “Codebook Varieties of Party Identity and
Organization (V-Party) V2”. Varieties of Democracy (V-Dem) Project, University of Gothenburg.
\url{https://doi.org/10.23696/vpartydsv2}

\section*{Tabla de control de versiones}
\begin{longtable}{@{}p{1.4cm}p{1.9cm}p{9.4cm}p{2.5cm}@{}}
\toprule
\textbf{Versión} & \textbf{Fecha} & \textbf{Cambios} & \textbf{Autor(es)}\\
\midrule\endhead
1.0 & Ago. 2026 & Primera versión: 8 dimensiones ideológicas transcritas desde la matriz de definición de
ejes; estructura documental alineada al formato V-Party. & Trujillo / Camargo\\
\addlinespace
1.1 & Oct. 2026 &
\textbf{Calibración tras la primera ronda de anotación (n = 207).}
(a) Dimensión Institucionalismo: se retira el marcador «formalidad burocrática y neutralidad», se
reescriben los marcadores para exigir función argumentativa, se añade una lista de \emph{no-marcadores} y
se redefinen los niveles 2 y 3 de la escala; se añade un ejemplo de Nivel 2.
(b) §2.1: se ratifican las reglas de agregación, empate y umbral mínimo, con criterio de desempate
determinístico asimétrico.
(c) §2.2 (nueva): procedimiento abreviado de anotación.
(d) §1.5: nueva instrucción general «el tema no es el encuadre» y prueba de sustracción.
(e) §1.6 y §1.5: se ratifican la unidad de codificación y el tratamiento de citas textuales.
(f) Apéndice D (nuevo): registro razonado de los cambios y evidencia que los motiva.
(g) Corrección de referencias cruzadas a §1.5.1 (inexistente) por §1.4.1.
& Trujillo / Camargo\\
\bottomrule
\end{longtable}

\caja{nuevo}{Cómo leer esta versión}{%
Todos los cambios respecto de v1.0 están justificados con evidencia en el \textbf{Apéndice D}. Las
decisiones que en v1.0 aparecían marcadas como «pendientes de ratificación por el equipo» quedan resueltas
en v1.1 y señaladas con el rótulo \textbf{Ratificada (v1.1)}. Ninguna de ellas modifica el Anteproyecto: la
arquitectura comprometida sigue siendo una única distribución \emph{softmax} sobre las ocho clases.}

\newpage
\tableofcontents
\newpage

% ============================ 1. NOTAS EXPLICATORIAS ============================
\section{Notas Explicatorias}

\subsection{Propósito de este Codebook}
Este codebook operacionaliza la identificación de la orientación ideológica dominante en noticias políticas
colombianas, en el marco del proyecto de grado “Desarrollo de clasificador multiclase para la identificación
de orientación ideológica en noticias políticas colombianas”. Define ocho (8) dimensiones ideológicas
medibles de forma independiente, cada una en una escala ordinal de 1 (Ausente) a 5 (Dominante), y
establece las reglas para derivar de ellas una única clase ideológica predominante por texto.

El codebook cumple dos funciones simultáneas en el pipeline del proyecto: (1) es el instrumento que usan
los dos investigadores para anotar manualmente el Gold Set (conjunto de calibración humana), y (2) es el
conjunto de reglas operacionales que se entrega como instrucción a un modelo de lenguaje de gran escala
(LLM) para generar el Silver Set (etiquetado masivo asistido), cuya consistencia se audita contra el consenso
humano mediante el Alfa de Krippendorff ($\alpha \geq 0.8$).

\caja{ratif}{Consecuencia operativa de la doble función}{%
Las dos funciones comparten un solo texto: cualquier cambio en este documento debe propagarse
simultáneamente al instrumento de anotación humana y a las instrucciones del juez LLM. Si solo se corrige
uno de los dos, el Silver Set deja de ser comparable con el Gold Set y la auditoría de concordancia pierde su
sentido, porque mediría la diferencia entre dos codebooks y no la fidelidad del etiquetado asistido.}

\subsection{Marco Teórico: Cuatro Ejes, Ocho Dimensiones}
Siguiendo la Teoría del Encuadre (\emph{framing}) de Entman, se asume que los medios no solo determinan qué
temas cubrir sino cómo enmarcarlos. Para operacionalizar este encuadre, la taxonomía se organiza en cuatro
ejes teóricos bipolares, cada uno compuesto por dos dimensiones que se codifican de forma independiente
(no mutuamente excluyente). La taxonomía se alinea conceptualmente con las variables e índices del V-Party
Dataset (V-Dem), adaptados a la agenda política colombiana.

\begin{longtable}{@{}p{0.6cm}p{4.1cm}p{3.6cm}p{6.6cm}@{}}
\toprule
\textbf{\#} & \textbf{Eje teórico} & \textbf{Dimensiones} & \textbf{Correspondencia V-Party}\\
\midrule\endhead
1 & Progresismo vs. Conservadurismo & Progresismo vs. Conservadurismo & Derechos civiles, diversidad de
género y secularismo (V-Party)\\
\addlinespace
2 & Populismo vs. Institucionalismo & Populismo vs. Institucionalismo & Retórica populista (anti-elitismo y
centralidad de la voluntad popular) frente al compromiso con el pluralismo (V-Party)\\
\addlinespace
3 & Soberanismo vs. Globalismo & Soberanismo vs. Globalismo & Nacionalismo económico frente a
internacionalismo y multilateralismo (marco teórico propio; sin variable V-Party directa)\\
\addlinespace
4 & Personalismo vs. Doctrinarismo & Personalismo vs. Doctrinarismo & Personalismo de partido (V-Party);
Doctrinarismo es contraparte definida por el proyecto\\
\bottomrule
\end{longtable}

\subsection{Notas de Precaución}

\caja{ratif}{Estado de las decisiones pendientes de v1.0 — Ratificadas (v1.1)}{%
La versión 1.0 identificó cuatro decisiones metodológicas sin resolver que debían fijarse por consenso de los
dos investigadores antes de iniciar la anotación del Gold Set: (1) la unidad exacta de codificación, (2) la
regla de derivación de la clase predominante, (3) el umbral de «sin encuadre dominante», y (4) el
tratamiento de citas textuales de terceros dentro de la noticia.\par\smallskip
\textbf{Las cuatro quedan resueltas en esta versión}, tras la primera ronda de anotación: (1) en §1.6,
(2) y (3) en §2.1, y (4) en §1.5. El Apéndice D documenta la evidencia que motivó cada decisión. La única
cuestión que permanece abierta —y que requeriría aprobación del director por modificar el Anteproyecto— es
si la categoría «Sin encuadre ideológico dominante» debe convertirse en una novena clase del clasificador;
v1.1 la trata como criterio de exclusión y \emph{no} como clase (§2.1, regla 3).}

Como en V-Party, se advierte que la orientación ideológica es un constructo intersubjetivo: el Gold Set no
captura una “verdad absoluta”, sino el consenso estructurado de los codificadores bajo las reglas de este
manual (Anteproyecto, sección 4.2). Por ello, toda ambigüedad debe documentarse en el campo de
comentarios (Apéndice C) en lugar de resolverse de forma silenciosa e inconsistente entre anotadores.

\subsection{Información de las Variables}

\subsubsection{Tipos de Variable}
Este codebook define dos tipos de variable, análogos a los tipos A/C de V-Dem:
\begin{itemize}
\item \textbf{Tipo H} (\emph{Human} — Gold Set): codificada de forma independiente por los dos
investigadores siguiendo estrictamente este manual. Las discrepancias se resuelven por consenso.
Constituye el estándar de verdad (\emph{Ground Truth}).
\item \textbf{Tipo LLM} (Silver Set): codificada por un modelo de lenguaje de gran escala instruido con este
mismo codebook. Sujeta a auditoría y curaduría manual sobre muestras aleatorias, validada contra el
consenso humano mediante Alfa de Krippendorff.
\item \textbf{Tipo D} (Índice derivado): variable calculada a partir de las 8 dimensiones anteriores mediante
una regla de agregación explícita (no observada directamente). Corresponde a la Clase Ideológica
Predominante (§2).
\end{itemize}

\subsubsection{Número de Variables}
\begin{longtable}{@{}p{7.2cm}p{3.2cm}p{3.0cm}@{}}
\toprule
\textbf{Sección} & \textbf{Tipo} & \textbf{\# Variables}\\
\midrule\endhead
Índice derivado (clase predominante) & D & 1\\
Eje 1 — Sociocultural & H, LLM & 2\\
Eje 2 — Gobernanza & H, LLM & 2\\
Eje 3 — Política exterior & H, LLM & 2\\
Eje 4 — Estructura de liderazgo & H, LLM & 2\\
\midrule
\textbf{Total} & & \textbf{9}\\
\bottomrule
\end{longtable}

\subsubsection{Convención de Nombres (Tags)}
Cada variable recibe un identificador corto, análogo a los tags \texttt{v2pa…} de V-Party, con la estructura:
prefijo \texttt{id\_} + abreviatura de 3–4 letras de la dimensión.

\begin{longtable}{@{}p{2.2cm}p{5.0cm}p{2.2cm}p{5.0cm}@{}}
\toprule
\textbf{Tag} & \textbf{Dimensión} & \textbf{Tag} & \textbf{Dimensión}\\
\midrule\endhead
\texttt{id\_prog} & Progresismo      & \texttt{id\_cons} & Conservadurismo\\
\texttt{id\_popu} & Populismo        & \texttt{id\_inst} & Institucionalismo\\
\texttt{id\_sob}  & Soberanismo      & \texttt{id\_glob} & Globalismo\\
\texttt{id\_pers} & Personalismo     & \texttt{id\_doct} & Doctrinarismo\\
\bottomrule
\end{longtable}
Variable derivada: \texttt{id\_pred} (Clase Ideológica Predominante).

\subsubsection{Campos de cada Ficha de Variable}
Cada una de las 8 dimensiones (§3) se documenta con la siguiente información, cuando aplica:
\begin{itemize}
\item \textbf{Eje teórico}: eje bipolar al que pertenece la dimensión.
\item \textbf{Pregunta}: definición operacional que responde la escala (equivalente al campo \emph{Question}
de V-Party).
\item \textbf{Marcadores léxicos y discursivos}: indicadores textuales, léxicos y retóricos que evidencian la
presencia de la dimensión (equivalente a \emph{Clarification}).
\item \textbf{No son marcadores}: rasgos que podrían confundirse con la dimensión pero que, por sí solos, no
la activan. Se documentan para las dimensiones donde la primera ronda de anotación reveló confusión
sistemática (ver Apéndice D).
\item \textbf{Escala}: niveles ordinales 1 a 5, cada uno con su descripción operacional (equivalente a
\emph{Responses}).
\item \textbf{Ejemplos}: casos ilustrativos tomados de la matriz de definición del proyecto y del corpus.
\item \textbf{Correspondencia V-Party}: variable o índice del V-Party Dataset del que se deriva
conceptualmente la dimensión (equivalente a \emph{Source}).
\item Tipo de variable, unidad de codificación, agregación entre codificadores y versión del codebook.
\end{itemize}

\subsection{Instrucciones Generales a los Codificadores}
\emph{(Adaptado del formato “Instructions to the coders” de V-Party, sección 3 del codebook original)}
\begin{itemize}
\item Lea el artículo completo antes de asignar cualquier puntaje. No codifique a partir del titular o del
primer párrafo únicamente.
\item Las ocho dimensiones se puntúan de forma independiente en una escala de 1 a 5. No son mutuamente
excluyentes: un mismo texto puede (y frecuentemente lo hará) obtener puntajes intermedios en varias
dimensiones simultáneamente.
\item El puntaje refleja \textbf{qué tanto estructura el argumento del texto} esa dimensión, no si el tema
aparece de forma incidental. Una sola mención aislada de un actor internacional no basta para un puntaje
alto en Globalismo; debe ser parte del argumento central.
\item Codifique el encuadre del medio, no la veracidad de los hechos relatados. Una noticia puede ser
factualmente correcta y aun así estar fuertemente enmarcada en una dimensión ideológica.
\item Registre cualquier caso ambiguo, límite o de difícil decisión en el campo de comentarios (Apéndice C),
incluyendo una breve justificación del puntaje asignado. Esto es insumo directo para iterar el codebook si el
Alfa de Krippendorff no alcanza el umbral de 0.8.
\end{itemize}

\caja{nuevo}{Nuevo en v1.1 — El tema no es el encuadre}{%
\textbf{El asunto del que trata una noticia no determina ninguno de los ocho puntajes.} Una nota sobre
pensiones puede puntuar alto en Conservadurismo, en Progresismo o en Institucionalismo según cómo la
enmarque; que hable de pensiones no aporta evidencia para ninguna de las tres.\par\smallskip
\textbf{Prueba de sustracción.} Antes de asignar un puntaje de 3 o más, pregúntese: \emph{si retiro del texto
la dimensión que estoy evaluando, ¿se cae el argumento o solo cambia el protagonista?} Si el argumento
sigue en pie, la dimensión está presente pero no es dominante, y el puntaje no debe pasar de 2.\par\smallskip
Esta regla importa especialmente en \textbf{Institucionalismo} (§3.2), donde el vocabulario de la
administración pública —entidades, resoluciones, trámites— aparece en casi cualquier noticia política
colombiana sin que ello constituya un encuadre ideológico. La primera ronda de anotación mostró que esta
confusión, no corregida, absorbe la mayoría del corpus (Apéndice D).}

\caja{ratif}{Ratificada (v1.1) — Tratamiento de citas textuales de terceros}{%
Las citas textuales de actores entrevistados dentro de la noticia \textbf{sí cuentan} para la codificación, en
tanto reflejan una decisión editorial de selección y espacio por parte del medio. Si el mismo artículo presenta
un contrapeso explícito (declaración de otro actor, desmentido, contexto que matiza la cita), el puntaje de
dominancia de esa dimensión \textbf{debe moderarse} en al menos un nivel.\par\smallskip
\emph{Justificación:} el objeto de medición declarado en el Anteproyecto es el encuadre del medio, y la
selección de a quién se cita y cuánto espacio se le concede es precisamente una operación de encuadre. La
regla de moderación evita el efecto contrario: que un medio que contrasta dos posturas quede codificado con
el puntaje de la más estridente de las dos.}

\subsection{Unidad de Codificación}

\caja{ratif}{Ratificada (v1.1) — Artículo periodístico completo}{%
La unidad de codificación es el \textbf{artículo periodístico completo} (titular + cuerpo), tal como lo extrae el
scraper, previo a cualquier segmentación (\emph{chunking}) usada en la Fase 2 de entrenamiento
(Anteproyecto, sección 6, Fase 2). Es la unidad que corresponde a la decisión editorial completa del
medio.\par\smallskip
\emph{Justificación:} es además la unidad que el pipeline ya implementa de extremo a extremo —el juez, la
construcción del dataset y la inferencia componen la misma entrada canónica (titular + cuerpo)—, de modo
que mantener la etiqueta a nivel de artículo evita una divergencia entre lo anotado y lo entrenado.\par\smallskip
La opción (b) de v1.0 —añadir una capa de anotación a nivel de \emph{chunk} para estudiar la dilución de
señal semántica declarada como limitación en el Anteproyecto (sección 4.2)— \textbf{no se descarta}, pero
queda explícitamente fuera del alcance de esta iteración: se documenta como trabajo futuro, no como
pendiente del Gold Set.}

\subsection{Protocolo de Concordancia entre Codificadores}
Siguiendo el Anteproyecto (sección 5.3, “Construcción del Corpus”), el proceso se estructura en dos fases
cuantificables mediante el Alfa de Krippendorff ($\alpha$):
\begin{itemize}
\item \textbf{Calibración humana (Gold Set)}: los dos investigadores anotan de forma independiente una
muestra inicial del corpus, estrictamente según este codebook. Se calcula $\alpha$ sobre las 8 dimensiones.
Si $\alpha \geq 0.8$, las discrepancias se resuelven por consenso y esa muestra pasa a ser el estándar de
verdad. Si $\alpha < 0.8$, \textbf{el codebook se itera} (se refinan definiciones, marcadores o niveles
ambiguos) y se repite la ronda de calibración.
\item \textbf{Etiquetado asistido (Silver Set)}: un LLM codifica el resto del corpus siguiendo este mismo
codebook. Se audita mediante muestras aleatorias, calculando $\alpha$ entre las etiquetas del LLM y el
consenso de los investigadores en esas muestras, para garantizar que el Silver Set no sufra degradación
conceptual respecto al Gold Set.
\end{itemize}

\caja{nuevo}{Nota de procedencia de v1.1}{%
Esta versión es, formalmente, una \textbf{iteración del codebook} en el sentido del protocolo anterior: la
primera ronda de anotación reveló un nivel ambiguo —el Nivel 3 de Institucionalismo— que concentraba el
81,6\,\% de las asignaciones. La iteración se ejecutó sobre la evidencia de esa ronda y antes de completar
el Gold Set, que es el momento previsto para hacerlo. El Apéndice D documenta el diagnóstico.}

\newpage
% ============================ 2. ÍNDICE DERIVADO ============================
\section{Índice Derivado}

\subsection{Clase Ideológica Predominante (D) — \texttt{id\_pred}}
\textbf{Pregunta:} ¿Cuál es la orientación ideológica dominante del texto, considerando las 8 dimensiones
codificadas?

\textbf{Fuente(s):} \texttt{id\_prog}, \texttt{id\_cons}, \texttt{id\_popu}, \texttt{id\_inst},
\texttt{id\_sob}, \texttt{id\_glob}, \texttt{id\_pers}, \texttt{id\_doct}

\textbf{Respuestas:} Categórica, una de las 8 dimensiones, o la categoría adicional
\textbf{“Sin encuadre ideológico dominante”}.

\caja{ratif}{Ratificada (v1.1) — Regla de agregación}{%
El Excel de definición asigna un puntaje 1–5 independiente a cada una de las 8 dimensiones, pero el
clasificador final necesita una única etiqueta por texto (\emph{softmax} de 8 clases, Anteproyecto §5.3). La
regla queda fijada así:\par\medskip
\textbf{(1) Regla general.} \texttt{id\_pred} = \emph{argmax} de los 8 puntajes (la dimensión con el puntaje
más alto).\par\medskip
\textbf{(2) Regla de empate.} Si dos o más dimensiones empatan en el puntaje máximo, el anotador decide
cuál estructura de forma más central el argumento del texto y registra la justificación en el campo de
comentarios (Apéndice C). Cuando no sea posible decidir, se aplica el siguiente criterio determinístico de
respaldo, \textbf{en orden}:\par\smallskip
\hspace{1em}\textbf{(a)} Si el empate involucra a \textbf{Institucionalismo} y a cualquier otra dimensión,
gana \textbf{la otra}.\par
\hspace{1em}\textbf{(b)} Si el empate persiste, gana la dimensión cuyos marcadores aparecen primero en el
titular y el primer tercio del texto.\par
\hspace{1em}\textbf{(c)} Si aún persiste, el caso se marca \texttt{empate\_irresoluble} en comentarios y se
excluye del Gold Set inicial, documentando su frecuencia en el reporte P1.3.\par\medskip
\emph{Justificación de la asimetría en (a):} Institucionalismo es la única de las ocho dimensiones cuyos
marcadores superficiales —léxico jurídico, entidades del Estado, descripción de trámites— aparecen en
prácticamente cualquier noticia política colombiana con independencia de su encuadre. Tiene, por tanto, una
tasa base de activación espuria estructuralmente más alta que las demás, confirmada empíricamente en la
primera ronda de anotación (Apéndice D). La asimetría \textbf{corrige un sesgo de medición conocido y
documentado}; no expresa una preferencia teórica por ninguna corriente ni altera la definición de la
dimensión.}

\caja{ratif}{Ratificada (v1.1) — Regla de umbral mínimo}{%
\textbf{(3)} Si el puntaje máximo entre las 8 dimensiones es \textbf{$\leq$ 2} —es decir, ninguna dimensión
supera el nivel «Leve»— el texto se etiqueta como \textbf{“Sin encuadre ideológico dominante”}.\par\medskip
\textbf{Tratamiento:} estos textos se \textbf{excluyen} del conjunto de entrenamiento y del conjunto de
prueba. No constituyen una novena clase. Su proporción sobre el corpus anotado se reporta como
característica del corpus en el producto P1.3.\par\medskip
\emph{Justificación:} el Anteproyecto compromete una única distribución \emph{softmax} de dimensión 8.
Convertir «Sin encuadre» en novena clase modificaría la arquitectura comprometida y requeriría aprobación
del director; excluir esos textos, en cambio, es un criterio de inclusión de datos y no toca la arquitectura.
Forzar una de las ocho clases sobre un texto sin encuadre es la alternativa peor: introduce ruido de etiqueta
en la clase que más se le parezca superficialmente —en la práctica, Institucionalismo— y el modelo aprende
esa confusión como si fuera señal.\par\medskip
\emph{Abierto para una versión futura:} reincorporar «Sin encuadre» como novena clase es defendible —haría
al clasificador utilizable sobre un corpus sin filtrar previamente— pero exige modificar el Anteproyecto y, por
tanto, el visto bueno del director. Queda registrado aquí como decisión diferida, no descartada.}

\textbf{Tipo de variable:} D (derivada; no se codifica directamente, se calcula a partir de las 8 variables
H/LLM)\\
\textbf{Versión del codebook:} 1.1

\subsection{Procedimiento Abreviado de Anotación \normalfont\small\color{gris}(nuevo en v1.1)}

El procedimiento completo descrito en §2.1 exige puntuar las ocho dimensiones y derivar de ellas el
\emph{argmax}: nueve decisiones por artículo. Sobre un Gold Set de más de mil artículos, ese costo es la
diferencia entre un conjunto anotado y uno a medio anotar. Se admite por ello un procedimiento abreviado:

\begin{enumerate}
\item El anotador identifica \textbf{directamente} la dimensión dominante, aplicando el criterio de
dominancia (§1.5) y la prueba de sustracción.
\item Asigna el puntaje 1–5 \textbf{únicamente a esa dimensión} (\texttt{id\_pred\_fuerza}).
\item Si ese puntaje es $\leq$ 2, la etiqueta pasa a ser “Sin encuadre ideológico dominante” (§2.1, regla 3).
\end{enumerate}

\caja{nuevo}{Equivalencia y límites del procedimiento abreviado}{%
\textbf{Es equivalente para la regla de umbral.} La regla 3 solo necesita saber si el \emph{máximo} de los
ocho puntajes supera 2. Si el anotador identificó correctamente la dimensión dominante, su puntaje
\emph{es} ese máximo, de modo que el umbral se aplica con un puntaje en lugar de ocho.\par\smallskip
\textbf{Lo que se pierde} es el perfil completo de las ocho dimensiones, que tiene dos usos: es señal
secundaria para el entrenamiento y es la base sobre la que se calcula el $\alpha$ de Krippendorff
\emph{por dimensión}.\par\smallskip
\textbf{Por tanto:} el procedimiento \textbf{completo} (ocho puntajes) es obligatorio en la muestra de
calibración sobre la que se calcula $\alpha$. El procedimiento \textbf{abreviado} se admite en el resto del
Gold Set. Cuál de los dos se usó debe quedar registrado por artículo en la columna \texttt{procedimiento} de
la plantilla (Apéndice B), porque de ello depende qué artículos entran en el cálculo de concordancia.}

\caja{pend}{Advertencia de implementación}{%
La interfaz de anotación debe capturar \texttt{id\_pred\_fuerza} además de \texttt{id\_pred}. Una interfaz que
solo ofrezca la clase dominante, sin puntaje, \textbf{hace inaplicable la regla de umbral}: obliga al anotador
a elegir una de las ocho clases para todo artículo, incluidos los que no tienen encuadre. Esa fue la causa
inmediata de la distribución documentada en el Apéndice D.}

\newpage
% ============================ 3. DIMENSIONES ============================
\section{Dimensiones Ideológicas}
\emph{Instrucciones a los codificadores} (ver también §1.5): evalúe cada una de las siguientes 8 dimensiones
de manera independiente sobre el mismo texto, antes de aplicar la regla de clase predominante de §2.

\subsection{Eje Sociocultural: Progresismo vs. Conservadurismo}

\dimtitulo{Progresismo}{id\_prog}
\textbf{Eje teórico:} Eje Sociocultural: Progresismo vs. Conservadurismo

\textbf{Pregunta:} Grado en que el texto impulsa transformaciones del \emph{statu quo} en favor de mayor
igualdad, justicia social, derechos de grupos históricamente marginados (mujeres, comunidades étnicas,
LGBTQ+, víctimas del conflicto, campesinado) o protección ambiental, usando estas causas como marco de
legitimación política.

\textbf{Marcadores léxicos y discursivos:}
\begin{itemize}
\item Derechos de minorías (LGBTQ+, indígenas, afrodescendientes, mujeres)
\item Justicia climática y medioambiental
\item Redistribución económica y justicia social
\item Laicismo y separación iglesia-Estado
\item Igualdad de género y feminismo
\item Reforma agraria y acceso a la tierra
\item Justicia transicional y memoria histórica
\item Paz como proceso transformador (no solo ausencia de conflicto)
\item Se visibilizan voces o demandas de grupos marginados como sujetos políticos
\item Se identifican estructuras o sistemas como fuente de desigualdad a transformar (sistema de salud,
modelo agrario, etc.)
\item Se critica el \emph{statu quo} como reproductor de inequidad y se proponen cambios estructurales
\end{itemize}

\textbf{Escala (Ordinal, 1–5):}
\begin{longtable}{@{}>{\centering\arraybackslash}p{1.1cm}p{14.3cm}@{}}
\toprule \textbf{Nivel} & \textbf{Descripción operacional}\\ \midrule\endhead
1 & No hay reivindicaciones de cambio social, igualdad o derechos; el texto es neutro o conservador.\\
\addlinespace
2 & Mención a temas sociales o grupos vulnerables, sin que sea el argumento central del texto.\\
\addlinespace
3 & El texto incluye varias menciones a justicia social, inclusión o reforma, combinadas con otros temas.\\
\addlinespace
4 & La narrativa está orientada explícitamente a transformar inequidades; los derechos o la justicia social
son ejes centrales.\\
\addlinespace
5 & Todo el texto es una proclama de transformación radical, exigiendo cambios estructurales y
reivindicando derechos de grupos oprimidos.\\
\bottomrule
\end{longtable}

\ejemplo{Nivel 5 (Dominante)}{La paz real y duradera es imposible sin una reforma agraria integral que
devuelva la tierra al campesinado y a las comunidades indígenas. Nuestro imperativo es transformar las
estructuras de exclusión histórica, garantizando los derechos plenos de las mujeres y diversidades, mientras
aceleramos una transición justa para dejar atrás el modelo económico extractivista que destruye la vida.}

\ejemplo{Nivel 1 (Ausente)}{El gremio empresarial hizo un llamado urgente a mantener la estabilidad
jurídica y a no modificar la actual legislación laboral. Los voceros del sector aseguraron que cualquier intento
de reforma estructural en este momento ahuyentará la inversión extranjera y pondrá en riesgo el libre
desarrollo de los mercados, que son el único motor real para generar riqueza.}

\metaficha{\texttt{v2paminor} (Minority rights), \texttt{v2palgbt} (LGBT social equality),
\texttt{v2pagender} (Gender equality), \texttt{v2pawelf} (Welfare, orientación universalista)}

\dimtitulo{Conservadurismo}{id\_cons}
\textbf{Eje teórico:} Eje Sociocultural: Progresismo vs. Conservadurismo

\textbf{Pregunta:} Grado en que el texto enfatiza valores tradicionales, orden social, seguridad, mano dura,
propiedad privada, familia tradicional, religión, gradualismo y escepticismo ante cambios rápidos
(especialmente en paz, tierras, drogas).

\textbf{Marcadores léxicos y discursivos:}
\begin{itemize}
\item Defensa de la propiedad privada y libre mercado
\item Valores familiares tradicionales
\item Religiosidad y moral cristiana como fundamento social
\item Autoridad, orden, seguridad y mano dura
\item Resistencia al cambio social (“así siempre ha sido”, tradiciones)
\item Defensa del \emph{statu quo} y estabilidad institucional
\item Se defiende el orden público, la autoridad estatal o el control del crimen como prioridades políticas
\item Conservadurismo colombiano apoyado en tres pilares: ortodoxia institucional/económica (propiedad
privada, libre empresa), exigencia de orden y autoridad (mano dura frente al crimen o la protesta, respaldo
irrestricto a la Fuerza Pública), y moral tradicional/religiosa (familia clásica como núcleo de la sociedad)
\item Se asocian posiciones progresistas o reformistas con riesgo, caos o pérdida de valores
\end{itemize}

\textbf{Escala (Ordinal, 1–5):}
\begin{longtable}{@{}>{\centering\arraybackslash}p{1.1cm}p{14.3cm}@{}}
\toprule \textbf{Nivel} & \textbf{Descripción operacional}\\ \midrule\endhead
1 & No hay apelación a la tradición, al orden ni a la autoridad; el texto puede ser neutral o proclive al
cambio.\\
\addlinespace
2 & Alguna referencia al orden o valores conservadores, pero sin que articule el enfoque del artículo.\\
\addlinespace
3 & Aparecen varias referencias a la estabilidad, la moral tradicional o la fuerza pública como elementos
relevantes, pero no lo son todo.\\
\addlinespace
4 & La defensa de la familia, la fe, el orden público y la autoridad es un tema central y constante; el cambio
se califica como amenaza.\\
\addlinespace
5 & Texto saturado de retórica de mano dura, valores tradicionales y rechazo al cambio social.\\
\bottomrule
\end{longtable}

\ejemplo{Nivel 5 (Dominante)}{No vamos a permitir que bajo el pretexto de un falso ‘cambio’ se destruyan
las instituciones, se premie a los delincuentes y se amenace la propiedad privada. Aquí lo que se necesita es
respaldar a nuestra Fuerza Pública con firmeza y autoridad para recuperar el orden, mientras defendemos la
familia como núcleo sagrado de la sociedad frente a ideologías que buscan destruir nuestros valores.}

\ejemplo{Nivel 1 (Ausente)}{La reforma radicada busca transformar de raíz el modelo vigente mediante la
redistribución equitativa de las tierras improductivas y la garantía plena de los derechos sexuales y
reproductivos. Es imperativo dejar atrás las estructuras anticuadas que han perpetuado la desigualdad para
construir una sociedad nueva y moderna.}

\metaficha{\texttt{v2parelig} (Religious principles), \texttt{v2paculsup} (Cultural superiority),
\texttt{v2palgbt} (inverso — oposición a igualdad LGBT), \texttt{v2pawelf} (orientación \emph{means-tested})}

\newpage
\subsection{Eje de Gobernanza y Discurso Estamental: Populismo vs. Institucionalismo}

\dimtitulo{Populismo}{id\_popu}
\textbf{Eje teórico:} Eje de Gobernanza y Discurso Estamental: Populismo vs. Institucionalismo

\textbf{Pregunta:} Grado en que el texto usa retórica que divide a la sociedad en “pueblo puro/virtuoso”
versus “élite corrupta”, apelando directamente a la voluntad popular contra las instituciones tradicionales.

\textbf{Marcadores léxicos y discursivos:}
\begin{itemize}
\item Dicotomía explícita pueblo vs. élite (“los de arriba”, “los de siempre”, “la oligarquía”, “el pueblo”)
\item Lenguaje anti-establecimiento y deslegitimación de instituciones existentes
\item Apelación emocional directa a las masas
\item Simplificación de problemas complejos en “buenos vs. malos”
\item Demonización de adversarios políticos
\item Uso de “nosotros” inclusivo contra “ellos” corrupto/privilegiado
\item Se apela a la “voluntad popular” o “el mandato del pueblo” para justificar acciones
\item Se invocan emociones de agravio, injusticia histórica o reivindicación colectiva
\end{itemize}

\textbf{Escala (Ordinal, 1–5):}
\begin{longtable}{@{}>{\centering\arraybackslash}p{1.1cm}p{14.3cm}@{}}
\toprule \textbf{Nivel} & \textbf{Descripción operacional}\\ \midrule\endhead
1 & Texto técnico, neutral o diplomático, sin apelación emocional ni dicotomías. No se detecta la división
pueblo-élite; los actores se presentan sin esa carga antagónica.\\
\addlinespace
2 & Alguna referencia a desigualdad pero sin antagonismo explícito.\\
\addlinespace
3 & La narrativa establece tensión entre sectores populares y élites, pero de forma descriptiva más que
movilizadora.\\
\addlinespace
4 & El texto articula explícitamente la división pueblo/élite y la usa como marco interpretativo central.\\
\addlinespace
5 & Todo el mensaje se articula en torno a esa oposición; el pueblo es exaltado como la única fuente de
virtud y la élite es demonizada como responsable de todos los males, con alta simplificación y afectividad.\\
\bottomrule
\end{longtable}

\ejemplo{Nivel 5 (Dominante)}{Este no es el proyecto de las oligarquías que por décadas se han robado y
repartido el país a sus espaldas. Este es el mandato del pueblo, de las familias trabajadoras y de ‘los nadies’
que hoy se levantan contra una élite corrupta, privilegiada y excluyente que solo ha legislado para sus
propios bolsillos desde sus clubes privados.}

\ejemplo{Nivel 1 (Ausente)}{El informe del Departamento Administrativo Nacional de Estadística (DANE)
señaló una contracción del 1.2\,\% en la demanda interna. Los analistas del sector financiero explicaron que
este fenómeno responde a ajustes macroeconómicos y a la estabilización de las tasas de interés por parte de
la junta directiva del Banco de la República.}

\metaficha{\texttt{v2paanteli} (Anti-elitism), \texttt{v2papeople} (People-centrism), \texttt{v2xpa\_popul}
(Índice de Populismo, tipo D)}

% ---------------- INSTITUCIONALISMO (REESCRITO EN v1.1) ----------------
\dimtitulo{Institucionalismo}{id\_inst}
\textbf{Eje teórico:} Eje de Gobernanza y Discurso Estamental: Populismo vs. Institucionalismo

\caja{nuevo}{Revisado en v1.1}{%
Esta ficha se reescribió tras la primera ronda de anotación, en la que absorbió el 81,6\,\% de las
asignaciones. Cambian los marcadores, se añade una lista de \textbf{no-marcadores} y se redefinen los
niveles 2 y 3 de la escala. La \textbf{definición teórica de la dimensión no cambia}: sigue siendo el
compromiso con el pluralismo y el procedimiento, anclado en \texttt{v2paplur}. Lo que cambia es el umbral de
evidencia exigido para afirmar que ese compromiso está presente. Ver Apéndice D.}

\textbf{Pregunta:} Grado en que el texto \textbf{argumenta} que las instituciones, los procesos legales
formales y el Estado de derecho deben ser los ejes de la acción política. La legitimidad se ancla en normas,
procedimientos y organismos formales, y no en voluntades individuales.

\caja{tinte}{Criterio de activación}{%
La dimensión se activa cuando el texto \textbf{defiende, reclama o evalúa} el procedimiento institucional —no
cuando simplemente \textbf{ocurre dentro} de él. La pregunta operativa es: \emph{¿el texto sostiene que el
procedimiento importa, o solo relata un hecho que sucedió a través de una institución?}}

\textbf{Marcadores léxicos y discursivos:}
\begin{itemize}
\item Se argumenta que la \textbf{validez} de una decisión política depende de que se haya seguido el
procedimiento debido
\item Se defiende explícitamente la separación de poderes, el debido proceso, la independencia judicial o la
de los órganos de control
\item Se señala una \textbf{irregularidad procedimental} como el problema central del asunto relatado
\item Se condenan las vías de hecho, los decretos por fuera del trámite legislativo o los saltos de la legalidad
\item Se apela al pluralismo o al respeto de «las reglas del juego» frente a la voluntad mayoritaria o la
urgencia política
\item Se evalúa a un actor político —a favor o en contra— \textbf{por su respeto a las reglas
institucionales}, y no por sus resultados, su carisma o su ideología
\item Léxico jurídico (decretos, leyes, artículos, sentencias) \textbf{cuando sostiene el argumento}, no
cuando solo describe el hecho relatado
\end{itemize}

\caja{pend}{No son marcadores de Institucionalismo \normalfont\small(nuevo en v1.1)}{%
Los siguientes rasgos, \textbf{por sí solos}, no otorgan puntaje por encima de 2:\par\smallskip
\textbf{(a) Que una entidad pública sea el sujeto de la noticia.} «La Alcaldía entregó», «la Corporación
acompañó al municipio», «el Congreso aprobó», «la Agencia analizará». Es el modo normal de narrar la
gestión pública en Colombia, no un encuadre ideológico.\par\smallskip
\textbf{(b) El tono neutro, formal o burocrático.} La neutralidad es la \textbf{ausencia} de encuadre: así se
define, de hecho, el Nivel 1 de las otras siete dimensiones de este codebook. Contarla \emph{a favor} de
Institucionalismo convertiría a esta dimensión en el destino por defecto de todo texto sin encuadre, que es
exactamente lo que ocurrió en la primera ronda (Apéndice D). \emph{Este marcador figuraba en v1.0 y se
retira en v1.1.}\par\smallskip
\textbf{(c) La mención de una norma, resolución o sentencia como dato del hecho relatado.} Citar la
Resolución 0450 porque es la noticia no es lo mismo que argumentar que la Resolución 0450 debía expedirse
siguiendo cierto trámite.\par\smallskip
\textbf{(d) Que el tema sea legislativo, judicial o administrativo.} El tema no es el encuadre (§1.5). Una
noticia sobre un debate en el Congreso puede ser Personalismo (si gira sobre la figura del ponente),
Populismo (si lo enmarca como pueblo contra élite) o Progresismo (si lo enmarca como conquista de
derechos).}

\textbf{Escala (Ordinal, 1–5):}
\begin{longtable}{@{}>{\centering\arraybackslash}p{1.1cm}p{14.3cm}@{}}
\toprule \textbf{Nivel} & \textbf{Descripción operacional}\\ \midrule\endhead
1 & No hay referencias a instituciones, leyes ni procedimientos formales; o se los menciona únicamente para
desestimarlos, ignorarlos o pasar por encima de ellos.\\
\addlinespace
2 & \textbf{(Revisado en v1.1)} Las instituciones aparecen como \textbf{sujeto gramatical} de la acción («el
Congreso aprobó», «la Corte determinó», «la Alcaldía anunció») o se cita una norma, pero el texto se limita
a \textbf{informar}: no argumenta que el procedimiento importe. La narrativa es informativa, no valorativa.
\emph{Este es el nivel por defecto de la nota administrativa, de gestión o de comunicado institucional.}\\
\addlinespace
3 & \textbf{(Revisado en v1.1)} El texto \textbf{valora explícitamente} algún proceso formal: respalda un
fallo, defiende un trámite, o señala una irregularidad procedimental como problema. Ese marco convive con
otros (pugna entre líderes, posición ideológica, reivindicación social) sin predominar sobre ellos.\\
\addlinespace
4 & El texto argumenta de forma repetida y central que la validez de una decisión política depende de que se
haya seguido el procedimiento. Si hay polémica, el foco está en si se respetaron las reglas, no en las
personas.\\
\addlinespace
5 & El texto defiende activamente el debido proceso, la separación de poderes o el pluralismo como valor
superior. Hay condena explícita de las vías de hecho, de los decretos por fuera del trámite legislativo o de
los saltos de la legalidad; la voluntad política individual se subordina al marco institucional.\\
\bottomrule
\end{longtable}

\ejemplo{Nivel 5 (Dominante)}{La Sala Plena de la Corte Constitucional declaró inexequible el artículo 4 de
la reforma, tras concluir que durante el trámite legislativo en el Congreso de la República se vulneró el
principio de consecutividad. En su providencia, el alto tribunal exhortó al Ejecutivo a expedir una nueva
reglamentación ceñida a los parámetros legales vigentes.}

\caja{nuevo}{Ejemplo — Nivel 2 \normalfont\small(nuevo en v1.1: el caso que más se confundió)}{%
\emph{La Agencia Nacional de Infraestructura analizará la posibilidad de reubicar el peaje del kilómetro 7. La
entidad informó que la solicitud fue radicada por la administración municipal y que la decisión se tomará una
vez se complete el estudio técnico contratado para tal fin.}\par\smallskip
\textbf{Por qué es Nivel 2 y no Nivel 3 ni superior:} hay una entidad del Estado como sujeto, un trámite
(«radicada») y un procedimiento («estudio técnico»). Pero el texto \textbf{no argumenta} que ese
procedimiento deba seguirse: lo relata porque es el hecho noticioso. Aplicando la prueba de sustracción
(§1.5), si se retira el encuadre institucional el argumento no se cae —sigue habiendo una noticia sobre un
peaje—, solo cambia quién la protagoniza.\par\smallskip
\textbf{Qué lo convertiría en Nivel 4 o 5:} que el texto sostuviera, por ejemplo, que la reubicación
\emph{no puede} decidirse por fuera del estudio técnico, o que denunciara que se pretende reubicar el peaje
saltándose el trámite. Ahí el procedimiento pasa de ser el escenario a ser el argumento.}

\ejemplo{Nivel 1 (Ausente)}{¡Aquí no vinimos a pedir permiso, vinimos a hacer historia! Yo les garantizo que
las cosas se van a hacer a mi manera, porque es lo que el pueblo en las calles me está exigiendo. No me
importan las trabas que nos pongan los mismos de siempre, mi voluntad es cumplirles hoy mismo.}

\metaficha{\texttt{v2paplur} (Political pluralism); inverso conceptual de \texttt{v2xpa\_antiplural} (Índice de
Anti-pluralismo, tipo D)}

\newpage
\subsection{Eje de Política Exterior: Soberanismo vs. Globalismo}

\dimtitulo{Soberanismo}{id\_sob}
\textbf{Eje teórico:} Eje de Política Exterior: Soberanismo vs. Globalismo

\textbf{Pregunta:} Grado en que el texto prioriza la autonomía, la integridad territorial y la
autodeterminación de Colombia frente a actores externos (organismos internacionales, potencias extranjeras,
empresas transnacionales, migrantes), presentando la soberanía nacional como valor político supremo.

\textbf{Marcadores léxicos y discursivos:}
\begin{itemize}
\item Defensa de fronteras y soberanía territorial
\item Proteccionismo económico (producción nacional, industria local)
\item Rechazo a injerencia extranjera (“no nos van a imponer”)
\item Nacionalismo y orgullo patrio
\item Énfasis en recursos naturales propios y su control
\item Se apela a identidad nacional, recursos naturales propios o fronteras como bienes a defender
\item Rechazo a injerencia externa en el contexto colombiano: en la izquierda $\rightarrow$ proteccionismo
agrario, crítica a la política antidrogas de EE.\,UU. y a los TLC; en la derecha $\rightarrow$ defensa
territorial en San Andrés y rechazo a organismos como la ONU o la CIDH frente al orden público
\end{itemize}

\textbf{Escala (Ordinal, 1–5):}
\begin{longtable}{@{}>{\centering\arraybackslash}p{1.1cm}p{14.3cm}@{}}
\toprule \textbf{Nivel} & \textbf{Descripción operacional}\\ \midrule\endhead
1 & No hay alusiones a la soberanía ni a injerencia externa; la interacción internacional se presenta como
cooperación natural.\\
\addlinespace
2 & Mención puntual a soberanía o actores externos, sin que articule el argumento principal del texto.\\
\addlinespace
3 & Se expresan varias inquietudes sobre presiones externas o se defiende la capacidad de decidir, sin ser el
tema central.\\
\addlinespace
4 & La protección de la soberanía y la crítica a intervenciones foráneas son parte central del discurso;
lenguaje cargado de “defensa nacional”.\\
\addlinespace
5 & El texto se estructura como una defensa frontal de la soberanía, denunciando amenazas externas
(embajadas, ONG, organismos multilaterales) y exigiendo decisiones sin ataduras foráneas.\\
\bottomrule
\end{longtable}

\ejemplo{Nivel 5 (Dominante)}{No podemos permitir que tribunales internacionales, ONG extranjeras o
potencias nos vengan a dictar cómo manejar nuestro orden interno o la política antidrogas. La soberanía de
Colombia se respeta; las decisiones sobre nuestro territorio y nuestras leyes las tomamos los colombianos,
sin arrodillarnos ante imposiciones externas.}

\ejemplo{Nivel 1 (Ausente)}{El Ministerio presentó la nueva hoja de ruta económica, diseñada en estricta
coordinación con las recomendaciones del Fondo Monetario Internacional y la OCDE. El objetivo es homologar
nuestros estándares normativos para integrarnos plenamente a los mercados globales y atraer inversión
extranjera.}

\metaficha{Sin variable V-Party dedicada. Eje de política exterior derivado del marco teórico propio del
proyecto (Anteproyecto, §5.3, eje 3).}

\dimtitulo{Globalismo}{id\_glob}
\textbf{Eje teórico:} Eje de Política Exterior: Soberanismo vs. Globalismo

\textbf{Pregunta:} Grado en que el texto valora la cooperación internacional, la adhesión a normas globales y
la interdependencia como herramientas para abordar problemas políticos.

\textbf{Marcadores léxicos y discursivos:}
\begin{itemize}
\item Mención de organismos internacionales (ONU, OCDE, FMI, BID, OMS)
\item Tratados de libre comercio y apertura económica
\item Cooperación internacional y ayuda multilateral
\item Estándares y normativas globales
\item Inversión extranjera como valor positivo
\item Integración regional (CAN, Alianza del Pacífico, Mercosur)
\item Se compara la situación colombiana con tendencias o normas de otros países o regiones
\item Se apela a compromisos de Colombia ante la comunidad internacional
\end{itemize}

\textbf{Escala (Ordinal, 1–5):}
\begin{longtable}{@{}>{\centering\arraybackslash}p{1.1cm}p{14.3cm}@{}}
\toprule \textbf{Nivel} & \textbf{Descripción operacional}\\ \midrule\endhead
1 & No hay mención de organismos internacionales, normas globales ni cooperación; el enfoque es puramente
doméstico.\\
\addlinespace
2 & Mención incidental a contexto internacional, sin que sea argumento ni eje narrativo.\\
\addlinespace
3 & Las referencias a la comunidad internacional o a tratados complementan el análisis, pero no son el eje.\\
\addlinespace
4 & La cooperación internacional, los estándares multilaterales o las alianzas globales son un pilar repetido
del discurso.\\
\addlinespace
5 & Todo el mensaje está basado en la necesidad de alinearse a la comunidad internacional, cumplir tratados
y recurrir a mecanismos globales; se usan expresiones como “compromiso internacional”, “estándares
globales”, “responsabilidad compartida”, y se descalifica el aislacionismo.\\
\bottomrule
\end{longtable}

\ejemplo{Nivel 5 (Dominante)}{El canciller reiteró que el camino para el desarrollo del país exige consolidar
nuestra participación en la OCDE y fortalecer los lazos comerciales mediante la Alianza del Pacífico. Además,
destacó que la cooperación técnica del Banco Interamericano de Desarrollo y el cumplimiento de la agenda
climática global de la ONU serán la brújula para atraer la inversión extranjera que necesitamos.}

\ejemplo{Nivel 1 (Ausente)}{La Alcaldía anunció que destinará recursos propios para financiar a las
asociaciones de recicladores del municipio. El programa, que busca formalizar el trabajo de estas
cooperativas locales, será operado directamente por la Secretaría de Inclusión Social y dependerá del
presupuesto aprobado por el Concejo Municipal para el próximo año.}

\metaficha{Sin variable V-Party dedicada. Eje de política exterior derivado del marco teórico propio del
proyecto (Anteproyecto, §5.3, eje 3).}

\newpage
\subsection{Eje de Estructura de Liderazgo: Personalismo vs. Doctrinarismo}

\dimtitulo{Personalismo}{id\_pers}
\textbf{Eje teórico:} Eje de Estructura de Liderazgo: Personalismo vs. Doctrinarismo

\textbf{Pregunta:} Grado en que el texto enfatiza a un líder individual (su carisma, voluntad, decisiones o
figura) como eje central de la política, por encima de partidos o instituciones.

\textbf{Marcadores léxicos y discursivos:}
\begin{itemize}
\item Deícticos de primera persona (“yo decidí”, “mi gobierno”, “lo logró”, “visión de…”, “el presidente
ordenó”, “su pueblo”)
\item Nombre del líder como sujeto reiterado de la acción política
\item Retórica carismática, emocional y directa
\item Promesas personales o lenguaje de lealtad (“yo les prometo”, “conmigo va a cambiar”)
\item Culto a la personalidad o heroización del líder
\item Ausencia de referencias a procesos institucionales
\end{itemize}

\textbf{Escala (Ordinal, 1–5):}
\begin{longtable}{@{}>{\centering\arraybackslash}p{1.1cm}p{14.3cm}@{}}
\toprule \textbf{Nivel} & \textbf{Descripción operacional}\\ \midrule\endhead
1 & No hay referencia a figuras individuales. Texto puramente institucional.\\
\addlinespace
2 & Se menciona un líder de forma ocasional pero como parte de un proceso institucional, sin centralidad
como referente necesario.\\
\addlinespace
3 & El líder es un actor relevante y se mencionan algunas de sus características, pero no monopoliza la
narrativa.\\
\addlinespace
4 & La personalidad del líder y sus atributos (carisma, trayectoria, autenticidad) estructuran buena parte del
argumento.\\
\addlinespace
5 & El texto depende casi por completo de la figura del líder, exaltando sus virtudes y presentándolo como
motor único de la acción política, sin mediación institucional.\\
\bottomrule
\end{longtable}

\ejemplo{Nivel 5 (Dominante)}{El mandatario asumió todo el crédito por los avances del proyecto,
asegurando desde la tarima: ‘Fui yo quien tomó la decisión valiente de enfrentar a las mafias cuando nadie
más quería hacerlo. Mi gobierno ha logrado lo que no se hizo en 20 años’. Durante su discurso, la figura del
dirigente opacó por completo el trámite legislativo, presentándose como el salvador indispensable del
proceso.}

\ejemplo{Nivel 1 (Ausente)}{El Ministerio de Hacienda expidió la Resolución 0450, mediante la cual se
reglamentan las nuevas tarifas de retención. El documento técnico fue producto de varias mesas de diálogo
interinstitucional y su implementación estará a cargo de las direcciones regionales de la DIAN durante el
próximo trimestre.}

\caja{tinte}{Nota de coordinación con §3.2 \normalfont\small(nueva en v1.1)}{%
El ejemplo de Nivel 1 anterior ilustra la ausencia de \emph{Personalismo}. No debe leerse como evidencia de
\emph{Institucionalismo} alto: según la escala revisada de §3.2, ese mismo texto puntúa \textbf{2} en
\texttt{id\_inst}, porque las entidades son sujeto de la acción pero el texto no argumenta sobre el valor del
procedimiento. Si ninguna otra dimensión supera 2, el artículo se clasifica como «Sin encuadre ideológico
dominante» (§2.1, regla 3). Un texto administrativo neutro no es, por serlo, un texto institucionalista.}

\metaficha{\texttt{v2paind} (Personalization of party)}

\dimtitulo{Doctrinarismo}{id\_doct}
\textbf{Eje teórico:} Eje de Estructura de Liderazgo: Personalismo vs. Doctrinarismo

\textbf{Pregunta:} Grado en que el texto presenta posiciones políticas como derivadas de un cuerpo ideológico
rígido y cerrado (de cualquier signo), al que se apela como autoridad o guía, mostrando resistencia a la
adaptación pragmática o al diálogo con posturas distintas.

\textbf{Marcadores léxicos y discursivos:}
\begin{itemize}
\item Uso frecuente de palabras “-ismos” (socialismo, neoliberalismo, marxismo, uribismo)
\item Citas a teóricos, manifiestos o doctrinas políticas
\item Se invocan principios o textos doctrinarios (constitucionales, religiosos, marxistas, libertarios, etc.)
como argumento de autoridad inapelable
\item Abstracción ideológica elevada, alejada de lo concreto
\item Dogmatismo: rechazo absoluto de posiciones intermedias
\item Marco teórico como lente único para interpretar la realidad
\item Lenguaje prescriptivo basado en doctrina (“según la teoría de…”)
\item El lenguaje es categórico, sin matices: “siempre”, “nunca”, “es imposible que”, “la única vía es”
\end{itemize}

\textbf{Escala (Ordinal, 1–5):}
\begin{longtable}{@{}>{\centering\arraybackslash}p{1.1cm}p{14.3cm}@{}}
\toprule \textbf{Nivel} & \textbf{Descripción operacional}\\ \midrule\endhead
1 & No hay presencia de doctrinas rígidas; el texto es pragmático o acepta matices.\\
\addlinespace
2 & Alguna referencia a principios ideológicos, pero sin presentarlos como inapelables o excluyentes, de
forma pragmática.\\
\addlinespace
3 & Se advierten referencias recurrentes a un ideario definido, pero se admiten otros enfoques.\\
\addlinespace
4 & El texto articula sus argumentos principalmente desde una posición ideológica rígida y descalifica
posturas alternativas.\\
\addlinespace
5 & El texto es un manifiesto doctrinal puro; todo se subordina a la pureza ideológica y cualquier concesión
es tratada como traición; lenguaje categórico (“principios innegociables”, “traición a los fundamentos”).\\
\bottomrule
\end{longtable}

\ejemplo{Nivel 5 (Dominante)}{La crisis estructural que atravesamos es la consecuencia ineludible del
modelo neoliberal y la acumulación por desposesión del capitalismo periférico. Siguiendo la teoría de la
dependencia, es evidente que cualquier reformismo tibio está condenado al fracaso; la única vía histórica es
desmantelar la hegemonía burguesa y avanzar hacia una transición socialista que recupere los medios de
producción.}

\ejemplo{Nivel 1 (Ausente)}{Para mitigar los trancones en la entrada sur de la ciudad, la Secretaría de
Movilidad anunció que habilitará un carril de contraflujo entre las 5:00 a.m. y las 8:00 a.m. Adicionalmente,
se destinarán 2.000 millones de pesos para tapar los huecos del corredor principal, con el fin de agilizar la
velocidad promedio de los vehículos de carga y transporte público.}

\metaficha{Sin variable V-Party dedicada. Concepto complementario definido por el equipo del proyecto como
contraparte de \texttt{v2paind} (Personalization of party).}

\newpage
% ============================ APÉNDICES ============================
\appendix
\section{Glosario}
\begin{longtable}{@{}p{4.3cm}p{11.1cm}@{}}
\toprule \textbf{Término} & \textbf{Definición}\\ \midrule\endhead
\emph{Chunking} (segmentación por bloques) & División de un texto extenso en fragmentos de longitud fija
para superar el límite de 512 tokens de las arquitecturas tipo BERT-encoder.\\
\addlinespace
Codebook & Conjunto de categorías, definiciones y criterios operacionales para clasificar/anotar datos de
forma consistente entre codificadores.\\
\addlinespace
\emph{Framing} (encuadre) & Selección, énfasis y omisión intencional de aspectos de un hecho por parte de un
medio, que produce narrativas divergentes sobre un mismo evento (Entman, 1993).\\
\addlinespace
Gold Set & Conjunto de datos anotado manualmente por los investigadores, auditado y validado, usado como
estándar de verdad (\emph{Ground Truth}).\\
\addlinespace
Silver Set & Conjunto de datos etiquetado por un LLM siguiendo el codebook, auditado mediante muestras
aleatorias contra el Gold Set.\\
\addlinespace
\emph{Ground Truth} & Información de referencia considerada correcta para efectos de entrenamiento y
evaluación del modelo.\\
\addlinespace
Alfa de Krippendorff ($\alpha$) & Métrica de concordancia entre codificadores:
$\alpha = 1 - (D_o / D_e)$, donde $D_o$ es el desacuerdo observado y $D_e$ el desacuerdo esperado por azar.
Umbral de aceptación en este proyecto: $\alpha \geq 0.8$.\\
\addlinespace
\emph{Softmax} & Función de activación de la capa de salida del clasificador que normaliza 8 puntuaciones
crudas en una distribución de probabilidad cuya suma es 1.0.\\
\addlinespace
Unidad de codificación & Fragmento de texto (artículo completo, párrafo o \emph{chunk}) sobre el cual se
asigna un puntaje. Ver §1.6.\\
\addlinespace
Clase ideológica predominante & Etiqueta categórica única derivada del \emph{argmax} de las 8 dimensiones
ideológicas puntuadas para un mismo texto (ver §2).\\
\addlinespace
Tasa base de activación espuria \emph{(v1.1)} & Frecuencia con que los marcadores superficiales de una
dimensión aparecen en textos que no tienen ese encuadre. Una tasa alta hace que la dimensión gane
\emph{argmax} por acumulación de indicios débiles y no por dominancia real. Ver Apéndice D.\\
\bottomrule
\end{longtable}

\section{Plantilla de Captura de Datos — Gold Set}
Estructura de columnas sugerida para la hoja de cálculo o base de datos donde se registre la anotación
manual. Cada fila corresponde a un artículo (unidad de codificación, §1.6).

\begin{longtable}{@{}p{3.5cm}p{9.1cm}p{2.6cm}@{}}
\toprule \textbf{Columna} & \textbf{Descripción} & \textbf{Tipo}\\ \midrule\endhead
\texttt{article\_id} & Identificador único del artículo (del pipeline de scraping) & Texto\\
\texttt{medio} & Nombre del medio de comunicación & Texto\\
\texttt{segmento\_medio} & Nacional / independiente / institucional / regional / gremial (Anteproyecto 4.1) &
Categórica\\
\texttt{fecha\_publicacion} & Fecha de publicación de la noticia & Fecha\\
\texttt{url} & Enlace original de la noticia & Texto\\
\addlinespace
\texttt{id\_prog} & Puntaje Progresismo & 1–5\\
\texttt{id\_cons} & Puntaje Conservadurismo & 1–5\\
\texttt{id\_popu} & Puntaje Populismo & 1–5\\
\texttt{id\_inst} & Puntaje Institucionalismo & 1–5\\
\texttt{id\_sob} & Puntaje Soberanismo & 1–5\\
\texttt{id\_glob} & Puntaje Globalismo & 1–5\\
\texttt{id\_pers} & Puntaje Personalismo & 1–5\\
\texttt{id\_doct} & Puntaje Doctrinarismo & 1–5\\
\addlinespace
\texttt{id\_pred} & Clase ideológica predominante (derivada, §2). Admite el valor
\texttt{sin\_encuadre}. \emph{(Ampliado en v1.1)} & Categórica\\
\texttt{id\_pred\_fuerza} & \emph{(Nuevo en v1.1)} Puntaje 1–5 de la dimensión señalada en
\texttt{id\_pred}. Obligatorio; es lo que permite aplicar la regla de umbral (§2.1, regla 3) & 1–5\\
\texttt{procedimiento} & \emph{(Nuevo en v1.1)} \texttt{completo} (8 puntajes) o \texttt{abreviado}
(§2.2). Determina si el artículo entra en el cálculo de $\alpha$ por dimensión & Categórica\\
\addlinespace
\texttt{codificador} & Investigador que asignó el puntaje & Texto\\
\texttt{comentarios} & Justificación de casos ambiguos o límite (ver Apéndice C) & Texto libre\\
\bottomrule
\end{longtable}

\section{Sección de Comentarios}
Análogo a las variables \texttt{v2padatacom} / \texttt{v2paccpar} del codebook V-Party. Cada codificador debe
dejar un comentario cuando:
\begin{itemize}
\item El puntaje asignado a alguna dimensión fue difícil de decidir entre dos niveles adyacentes.
\item El texto activó una dimensión no cubierta claramente por los marcadores léxicos listados.
\item Se aplicó la regla de empate o la regla de umbral mínimo de §2.1.
\item \emph{(Nuevo en v1.1)} Se recurrió al criterio determinístico de respaldo de §2.1, regla 2, indicando
cuál de los literales (a), (b) o (c) se aplicó.
\item Se considera que el codebook debería ajustarse (nueva definición, nuevo marcador, ejemplo adicional).
\end{itemize}
Estos comentarios no forman parte del dataset de entrenamiento, pero son el insumo principal para las
iteraciones de versión de este codebook (ver Tabla de control de versiones). El Apéndice D es un ejemplo de
ese uso.

\newpage
\section{Registro Razonado de Cambios (v1.0 $\rightarrow$ v1.1)}
\emph{Apéndice nuevo en v1.1. Documenta la evidencia que motivó cada cambio, en cumplimiento del
protocolo de iteración del codebook (§1.7).}

\subsection{Evidencia: primera ronda de anotación}
Sobre el Gold Set en construcción se completaron \textbf{207 artículos} anotados por un codificador con la
versión 1.0 del codebook, usando una interfaz que capturaba únicamente la clase dominante
(\texttt{id\_pred}), sin los ocho puntajes. La distribución resultante fue:

\begin{longtable}{@{}p{5.2cm}>{\centering\arraybackslash}p{2.6cm}>{\centering\arraybackslash}p{2.6cm}@{}}
\toprule \textbf{Clase asignada} & \textbf{n} & \textbf{\%}\\ \midrule\endhead
Institucionalismo & 169 & 81,6\,\%\\
Personalismo      &  10 &  4,8\,\%\\
Globalismo        &   9 &  4,3\,\%\\
Progresismo       &   9 &  4,3\,\%\\
Soberanismo       &   3 &  1,4\,\%\\
Conservadurismo   &   3 &  1,4\,\%\\
Populismo         &   2 &  1,0\,\%\\
Doctrinarismo     &   2 &  1,0\,\%\\
\midrule
\textbf{Total}    & \textbf{207} & \textbf{100\,\%}\\
\bottomrule
\end{longtable}

Para descartar que el sesgo proviniera de la composición del corpus —por ejemplo, de un exceso de fuentes
institucionales— se desagregó por categoría de medio. Institucionalismo domina en \textbf{todas} las
categorías, incluidas las de opinión y las independientes, que son las que por diseño deberían concentrar el
encuadre explícito:

\begin{longtable}{@{}p{4.4cm}>{\centering\arraybackslash}p{3.2cm}>{\centering\arraybackslash}p{3.0cm}@{}}
\toprule \textbf{Categoría de fuente} & \textbf{Institucionalismo} & \textbf{\%}\\ \midrule\endhead
Regional        & 77 / 87 & 88,5\,\%\\
Gremial         & 20 / 24 & 83,3\,\%\\
Opinión (\emph{think tanks}) & 16 / 20 & 80,0\,\%\\
Institucional   & 33 / 43 & 76,7\,\%\\
Nacional        & 11 / 15 & 73,3\,\%\\
Independiente   &  9 / 15 & 60,0\,\%\\
Judicial        &  3 /  3 & 100\,\%\\
\bottomrule
\end{longtable}

\subsection{Diagnóstico: tres causas convergentes}

\textbf{(1) Asimetría en el tratamiento de la neutralidad.} Las otras siete dimensiones usan la palabra
«neutral» para definir su \textbf{Nivel 1}, es decir, la \emph{ausencia} de la dimensión. Institucionalismo, en
cambio, listaba «formalidad burocrática y neutralidad» entre sus \textbf{marcadores a favor}. Un texto
periodístico neutro puntuaba por tanto 1 en siete dimensiones y subía en la octava, y el \emph{argmax} no
tenía escapatoria.

\textbf{(2) Umbral de evidencia demasiado bajo en el Nivel 3.} La redacción de v1.0 —«las instituciones o los
procedimientos legales articulan parte importante del argumento o la narrativa»— se satisface con la
presencia de una entidad pública como sujeto de la noticia. Esa es la forma normal de redactar casi cualquier
nota de gestión pública en Colombia, de modo que el Nivel 3 describía al periodismo político ordinario en
lugar de un encuadre ideológico.

\textbf{(3) Ausencia de válvula de escape en el instrumento.} La regla de umbral mínimo de §2.1 —que
asigna «Sin encuadre ideológico dominante» cuando ningún puntaje supera 2— existía en v1.0 pero estaba
marcada como pendiente de ratificación y \textbf{no se implementó en la interfaz de anotación}, que además
no capturaba los puntajes 1–5 necesarios para evaluarla. El codificador quedó obligado a elegir una de las
ocho clases para \emph{todo} artículo, incluidos los que no tenían encuadre. El excedente se acumuló en la
clase cuyos marcadores superficiales son más frecuentes.

Las tres causas apuntan en la misma dirección y se refuerzan entre sí: la primera y la segunda inflan el
puntaje de Institucionalismo en textos sin encuadre, y la tercera impide registrarlos como lo que son.

\subsection{Cambios derivados}
\begin{longtable}{@{}p{1.0cm}p{7.6cm}p{6.6cm}@{}}
\toprule \textbf{Causa} & \textbf{Cambio} & \textbf{Sección}\\ \midrule\endhead
(1) & Se retira «formalidad burocrática y neutralidad» de los marcadores de Institucionalismo. &
§3.2, marcadores\\
\addlinespace
(1) & Se añade una lista explícita de \emph{no-marcadores}, con los cuatro rasgos que no activan la dimensión
por sí solos. & §3.2, caja «No son marcadores»\\
\addlinespace
(1,\,2) & Se reescriben los marcadores para exigir función argumentativa («se argumenta que la validez
depende del procedimiento») en lugar de presencia léxica. & §3.2, marcadores\\
\addlinespace
(2) & Nivel 2 redefinido: institución como sujeto gramatical, narrativa informativa y no valorativa. Es el nivel
por defecto de la nota administrativa. & §3.2, escala\\
\addlinespace
(2) & Nivel 3 redefinido: exige valoración explícita de un proceso formal. & §3.2, escala\\
\addlinespace
(2) & Se añade un \textbf{Ejemplo de Nivel 2} con su explicación contrastiva, por ser el nivel que concentró
la confusión. & §3.2, ejemplos\\
\addlinespace
(2) & Instrucción general nueva: «el tema no es el encuadre» y prueba de sustracción, aplicable a las ocho
dimensiones. & §1.5\\
\addlinespace
(3) & Se ratifica la regla de umbral mínimo y se fija su tratamiento: exclusión, no novena clase. & §2.1,
regla 3\\
\addlinespace
(3) & Se ratifica la regla de empate con criterio determinístico asimétrico. & §2.1, regla 2\\
\addlinespace
(3) & Se define el procedimiento abreviado de anotación, que vuelve aplicable la regla de umbral con un solo
puntaje en lugar de ocho. & §2.2\\
\addlinespace
(3) & La plantilla incorpora \texttt{id\_pred\_fuerza} y \texttt{procedimiento}; \texttt{id\_pred} admite
\texttt{sin\_encuadre}. & Apéndice B\\
\bottomrule
\end{longtable}

\subsection{Qué no se cambió, y por qué}
\begin{itemize}
\item \textbf{Las ocho clases y los cuatro ejes.} Son el contrato del Anteproyecto aprobado y el anclaje con
el V-Party Dataset. v1.1 no añade, elimina ni renombra ninguna dimensión.
\item \textbf{La definición teórica de Institucionalismo.} Sigue siendo el compromiso con el pluralismo y el
procedimiento, anclada en \texttt{v2paplur}. Lo que v1.1 modifica es el \emph{umbral de evidencia} exigido
para afirmar que ese compromiso está presente en un texto, que es un asunto de operacionalización y no de
teoría.
\item \textbf{La arquitectura del clasificador.} Una única distribución \emph{softmax} de dimensión 8 que
suma 1,0, tal como está comprometida.
\item \textbf{El protocolo de concordancia.} El umbral $\alpha \geq 0.8$ y el procedimiento de doble
anotación independiente permanecen como están definidos en §1.7.
\end{itemize}

\subsection{Verificación prevista}
El efecto de esta revisión debe medirse antes de etiquetar el resto del Gold Set y antes de generar el Silver
Set:

\begin{enumerate}
\item \textbf{Re-anotar los 169 artículos} asignados a Institucionalismo en la primera ronda, con la escala
revisada y la válvula de escape habilitada. Los 38 restantes no están en disputa y no requieren revisión.
\item \textbf{Criterio de éxito esperado:} Institucionalismo debe seguir siendo la clase mayoritaria —el
corpus contiene una proporción alta de fuentes institucionales y comunicados oficiales, y esa composición es
real—, pero su participación debería situarse en el orden del 35–45\,\% y aparecer una proporción no trivial
de «Sin encuadre ideológico dominante». Una distribución uniforme entre las ocho clases \textbf{no} es el
objetivo y sería, de hecho, indicio de sobrecorrección.
\item \textbf{Propagar al juez LLM.} La definición revisada debe aplicarse también al codebook operacional
que instruye al modelo del Silver Set antes de ejecutarlo. Si el Silver Set se genera con la versión antigua, el
sesgo se reproduce a escala de todo el corpus y la auditoría de concordancia contra un Gold Set corregido lo
registraría como desacuerdo del juez, no como lo que sería: dos codebooks distintos (§1.1).
\end{enumerate}

\newpage
\section*{Referencias}
\addcontentsline{toc}{section}{Referencias}
\small
Entman, R. M. (1993). Framing: Toward Clarification of a Fractured Paradigm. \emph{Journal of
Communication}, 43(4), 51–58.

Hayes, A. F., \& Krippendorff, K. (2007). Answering the Call for a Standard Reliability Measure for Coding
Data. \emph{Communication Methods and Measures}, 1(1), 77–89.

Lindberg, S. I., Düpont, N., Higashijima, M., Kavasoglu, Y. B., Marquardt, K. L., Bernhard, M., Döring, H.,
Hicken, A., Laebens, M., Medzihorsky, J., Neundorf, A., Reuter, O. J., Ruth–Lovell, S., Weghorst, K. R.,
Wiesehomeier, N., Wright, J., Alizada, N., Bederke, P., Gastaldi, L., Grahn, S., Hindle, G., Ilchenko, N., von
Römer, J., Wilson, S., Pemstein, D., \& Seim, B. (2022). \emph{Codebook Varieties of Party Identity and
Organization (V-Party) V2}. Varieties of Democracy (V-Dem) Project.
\url{https://doi.org/10.23696/vpartydsv2}

Trujillo Riascos, K. S., \& Camargo Chacón, J. E. (2026). \emph{Desarrollo de clasificador multiclase para la
identificación de orientación ideológica en noticias políticas colombianas} [Anteproyecto de trabajo de
grado]. Universidad del Valle.

\end{document}
